

\setcounter{chapter}{7}

\chapter{欧几里得整环、主理想整环与唯一分解整环}

有若干类比一般环具有更多代数结构的环类。本章考虑的是具有除法算法的环（欧几里得整环）、每个理想都是主理想的环（主理想整环）以及元素可分解为素元的环（唯一分解整环）。这类环的主要例子是整数环 \(\mathbb{Z}\) 和系数在某个域 \(F\) 中的多项式环 \(F[x]\). 我们在这里把预备知识一章中关于整数 \(\mathbb{Z}\) 陈述的所有定理证明为对更一般环成立的结论的特殊情形。这些结果将在下一章应用于环 \(F[x]\) 的特殊情形。

本章中所有环都交换。

\section{欧几里得整环}

我们首先定义整环 \(R\) 上\textbf{范数}的概念。它本质上是 \(R\) 中“大小”的一种度量。

\begin{definition}
任何满足 \(N(0) = 0\) 的函数 \(N : R \to \mathbb{Z}^+ \cup \{0\}\) 称为整环 \(R\) 上的\textbf{范数}。若对 \(a \neq 0\) 有 \(N(a) > 0\)，则称 \(N\) 是\textbf{正范数}。
\end{definition}

我们注意到这个范数概念相当弱，同一个整环 \(R\) 可能拥有若干不同的范数。

\begin{definition}
整环 \(R\) 称为\textbf{欧几里得整环}（或具有\textbf{除法算法}），若存在 \(R\) 上的范数 \(N\)，使得对 \(R\) 的任意两个元素 \(a, b\)（\(b \neq 0\)），存在 \(R\) 中的元素 \(q, r\) 使
\[
a = qb + r, \quad \text{其中 } r = 0 \text{ 或 } N(r) < N(b).
\]
元素 \(q\) 称为除法的\textbf{商}，元素 \(r\) 称为\textbf{余数}。
\end{definition}

整环 \(R\) 上除法算法的存在的重要性在于它允许对 \(R\) 的两个元素 \(a, b\) 使用\textbf{欧几里得算法}：通过相继的“除法”（这些实际上是 \(R\) 的分式域中的除法）我们可以写出
\[
\begin{aligned}
a &= q_0b + r_0 \quad (0)\\
b &= q_1r_0 + r_1 \quad (1)\\
r_0 &= q_2r_1 + r_2 \quad (2)\\
&\ \ \vdots \\
r_{n-2} &= q_nr_{n-1} + r_n \quad (n)\\
r_{n-1} &= q_{n+1}r_n \quad (n+1)
\end{aligned}
\]
其中 \(r_n\) 是最后一个非零余数。这样的 \(r_n\) 存在，因为若余数非零，则 \(N(b) > N(r_0) > N(r_1) > \cdots > N(r_n)\) 是非负整数的递降序列，而这样的序列不能无限继续。还要注意这些元素不保证唯一。

\begin{example}
（0）域是欧几里得整环的平凡例子，任何范数都满足定义条件（例如对所有 \(a\) 取 \(N(a) = 0\)）。这是因为对每个满足 \(b \neq 0\) 的 \(a, b\) 有 \(a = qb + 0\)，其中 \(q = ab^{-1}\).

（1）整数 \(\mathbb{Z}\) 是欧几里得整环，范数由 \(N(a) = |a|\)（通常的绝对值）给出。\(\mathbb{Z}\) 中除法算法的存在（初等算术中熟悉的“长除法”）验证如下。设 \(a, b\) 是两个非零整数，先假设 \(b > 0\). 半开区间 \([nb, (n+1)b)\)（\(n \in \mathbb{Z}\)）划分实直线，故 \(a\) 在其中一个中，比如 \(a \in [kb, (k+1)b)\). 取 \(q = k\) 则 \(a - qb = r \in [0, |b|)\)，正合所需。若 \(b < 0\)（故 \(-b > 0\)），由刚才所证存在整数 \(q\) 使 \(a = q(-b) + r\)，其中 \(r = 0\) 或 \(|r| < |-b|\)；则 \(a = (-q)b + r\) 满足对 \(a, b\) 的除法算法要求。这个论证可以通过对 \(|a|\) 归纳变得更正式。

注意若 \(a\) 不是 \(b\) 的倍数，则 \(q, r\) 总有两种可能：上面的证明总产生正的余数 \(r\). 例如若 \(b > 0\) 且 \(q, r\) 如上且 \(r > 0\)，则 \(a = q'b + r'\)（其中 \(q' = q+1\)、\(r' = r-b\)）也满足应用于 \(a, b\) 的除法算法条件。因此 \(5 = 2\cdot 2 + 1 = 3\cdot 2 - 1\) 是把 \(\mathbb{Z}\) 中除法算法应用于 \(a = 5\)、\(b = 2\) 的两种方式。若要求余数非负，则商与余数唯一。

（2）若 \(F\) 是域，则多项式环 \(F[x]\) 是欧几里得整环，范数由 \(N(p(x)) = p(x)\) 的次数给出。多项式的除法算法就是多项式的“长除法”，对实系数多项式可能已经熟悉。其证明与 \(\mathbb{Z}\) 的情形非常相似，将在下一章给出（不过对多项式会证明商与余数唯一）。为使多项式环成为欧几里得整环，系数必须来自域，因为除法算法最终依赖于能够除任意非零系数。我们将在 8.2 节证明若 \(R\) 不是域则 \(R[x]\) 不是欧几里得整环。

（3）7.1 节的二次整数环 \(\mathcal{O}\) 是带有由域范数的绝对值定义的范数的整环（以保证取值为非负；当 \(D < 0\) 时域范数本身就是范数），但一般地 \(\mathcal{O}\) 对这个范数（或任何其他范数）不欧几里得。然而\textbf{高斯整数} \(\mathbb{Z}[i]\)（其中 \(D = -1\)）对范数 \(N(a+bi) = a^2+b^2\) 是欧几里得整环，我们现在证明这一点（另见 8.3 节末）。

设 \(\alpha = a+bi\)、\(\beta = c+di\) 是 \(\mathbb{Z}[i]\) 的两个元素且 \(\beta \neq 0\). 则在域 \(\mathbb{Q}(i)\) 中有 \(\frac{\alpha}{\beta} = r + si\)，其中 \(r = \frac{ac+bd}{c^2+d^2}\)、\(s = \frac{bc-ad}{c^2+d^2}\) 是有理数。设 \(p\) 是最接近有理数 \(r\) 的整数、\(q\) 是最接近有理数 \(s\) 的整数，于是 \(|r-p|\) 和 \(|s-q|\) 都不超过 \(\frac{1}{2}\). 一旦我们证明
\[
\alpha = (p+qi)\beta + \gamma, \quad \text{其中 } \gamma \in \mathbb{Z}[i] \text{ 且 } N(\gamma) \leq \tfrac{1}{2}N(\beta),
\]
除法算法就立即推出，而这比所需的更强。设 \(\delta = (r-p) + (s-q)i\)，令 \(\gamma = \beta\delta\). 则 \(\gamma = \alpha - (p+qi)\beta\)，故 \(\gamma \in \mathbb{Z}[i]\) 是高斯整数且 \(\alpha = (p+qi)\beta + \gamma\). 由于 \(N(\delta) = (r-p)^2 + (s-q)^2\) 至多为 \(\frac{1}{4} + \frac{1}{4} = \frac{1}{2}\)，范数 \(N\) 的乘性蕴含 \(N(\gamma) = N(\delta)N(\beta) \leq \frac{1}{2}N(\beta)\)，如所断言。

注意这个算法相当具体，因为商 \(p+qi\) 可以快速由有理数 \(r, s\) 确定，随后余数 \(\gamma = \alpha - (p+qi)\beta\) 容易计算。还要注意商不必唯一：若 \(r\)（或 \(s\)）是奇数的半数，则 \(p\)（或分别地 \(q\)）有两种选择。

这个证明 \(\mathbb{Z}[i]\) 是欧几里得整环的论证也可以用来证明 \(\mathcal{O}\) 对 \(D = -2, -3, -7, -11\) 是欧几里得整环（关于 7.1 节定义的域范数）（参见习题）。我们不久将看到 \(\mathbb{Z}[\sqrt{-5}]\) 对任何范数都不欧几里得，而 \(\mathbb{Z}[(1+\sqrt{-19})/2]\) 对任何范数都不是欧几里得整环的证明出现在本节末。

（4）回忆（参见 7.1 节习题 26）离散赋值环如下得到。设 \(K\) 是域。\(K\) 上的\textbf{离散赋值}是满足下列条件的函数 \(v : K^\times \to \mathbb{Z}\)：

（i）\(v(ab) = v(a) + v(b)\)（即 \(v\) 是从 \(K\) 的非零元素的乘法群到 \(\mathbb{Z}\) 的同态）；

（ii）\(v\) 是满射；

（iii）对所有满足 \(x+y \neq 0\) 的 \(x, y \in K^\times\) 有 \(v(x+y) \geq \min\{v(x), v(y)\}\).

集合 \(\{x \in K^\times \mid v(x) \geq 0\} \cup \{0\}\) 是 \(K\) 的子环，称为 \(v\) 的\textbf{赋值环}。整环 \(R\) 称为\textbf{离散赋值环}，若存在其分式域上的赋值 \(v\) 使 \(R\) 是 \(v\) 的赋值环。

例如分母与固定素数 \(p \in \mathbb{Z}\) 互素的所有有理数构成的环 \(R\) 是含于 \(\mathbb{Q}\) 的离散赋值环。

容易看出离散赋值环对由 \(N(0) = 0\)、在 \(R\) 的非零元素上 \(N = v\) 定义的范数是欧几里得整环。这是因为对满足 \(b \neq 0\) 的 \(a, b \in R\)：

（a）若 \(N(a) < N(b)\)，则 \(a = 0\cdot b + a\)；

（b）若 \(N(a) \geq N(b)\)，则由离散赋值的性质（i）可知 \(q = ab^{-1} \in R\)，故 \(a = qb + 0\).
\end{example}

整环 \(R\) 上除法算法的第一个推论是它迫使 \(R\) 的每个理想都是主理想。

\begin{proposition}\label{pro:8.1}
欧几里得整环中的每个理想都是主理想。更确切地说，若 \(I\) 是欧几里得整环 \(R\) 的任意非零理想，则 \(I = (d)\)，其中 \(d\) 是 \(I\) 中任意有最小范数的非零元素。
\end{proposition}

\begin{proof}
若 \(I\) 是零理想，无需证明。否则设 \(d\) 是 \(I\) 中任意有最小范数的非零元素（这样的 \(d\) 存在，因为由 \(\mathbb{Z}\) 的良序性，集合 \(\{N(a) \mid a \in I\}\) 有最小元）。显然 \((d) \subseteq I\)，因为 \(d\) 是 \(I\) 的元素。为证反向包含，设 \(a\) 是 \(I\) 的任意元素，用除法算法写 \(a = qd + r\)，其中 \(r = 0\) 或 \(N(r) < N(d)\). 则 \(r = a - qd\)，而 \(a\) 和 \(qd\) 都在 \(I\) 中，故 \(r\) 也是 \(I\) 的元素。由 \(d\) 的范数的极小性，我们看到 \(r\) 必为 0. 于是 \(a = qd \in (d)\)，表明 \(I = (d)\).
\end{proof}

命题 \ref{pro:8.1} 表明 \(\mathbb{Z}\) 的每个理想都是主理想。\(\mathbb{Z}\) 的这个基本性质先前（在 7.3 节）是用 \(\mathbb{Z}\) 的（加法）群结构、通过 2.3 节循环群子群的分类确定的。注意这实际上是同一个证明，因为 2.3 节的结果最终依赖于 \(\mathbb{Z}\) 中的欧几里得算法。

命题 \ref{pro:8.1} 也可以用来证明某些整环 \(R\) 不是（对任何范数的）欧几里得整环，方法是证明 \(R\) 中存在非主理想。

\begin{example}
（1）设 \(R = \mathbb{Z}[x]\)：由于理想 \((2, x)\) 不是主理想（参见 7.4 节开头的例 3），可知整系数多项式环 \(\mathbb{Z}[x]\) 不是（对任何范数的）欧几里得整环，尽管有理系数多项式环 \(\mathbb{Q}[x]\) 是欧几里得整环。

（2）设 \(R\) 是二次整数环 \(\mathbb{Z}[\sqrt{-5}]\)，设 \(N\) 是相应的域范数 \(N(a+b\sqrt{-5}) = a^2+5b^2\)，考虑由 3 和 \(2+\sqrt{-5}\) 生成的理想 \(I = (3, 2+\sqrt{-5})\). 假设 \(I = (\alpha + \beta\sqrt{-5})\)（\(\alpha, \beta \in \mathbb{Z}\)）是主的，即对某些 \(\alpha', \beta' \in R\) 有 \(3 = \alpha'(\alpha+\beta\sqrt{-5})\) 和 \(2+\sqrt{-5} = \beta'(\alpha+\beta\sqrt{-5})\). 在第一个等式中取范数得 \(9 = N(\alpha')(\alpha^2+5\beta^2)\)，由于 \(\alpha^2+5\beta^2\) 是正整数，它必为 1、3 或 9. 若这个值是 9，则 \(N(\alpha') = 1\) 且 \(\alpha' = \pm 1\)，故 \(\alpha+\beta\sqrt{-5} = \pm 3\)，由第二个等式这不可能，因为 \(2+\sqrt{-5}\) 的系数不被 3 整除。这个值不能是 3，因为 \(a^2+5b^2 = 3\) 无整数解。若这个值是 1，则 \(\alpha+\beta\sqrt{-5} = \pm 1\)，理想 \(I\) 将是整个环 \(R\). 但那样 1 就是 \(I\) 的元素，故对某些 \(\gamma, \delta \in R\) 有 \(3\gamma + (2+\sqrt{-5})\delta = 1\). 两边乘以 \(2-\sqrt{-5}\) 将推出 \(2-\sqrt{-5}\) 是 3 在 \(R\) 中的倍数，矛盾。由此推出 \(I\) 不是主理想，故 \(R\) 不是（对任何范数的）欧几里得整环。
\end{example}

\(\mathbb{Z}\) 中欧几里得算法的一个基本推论是它产生两个非零元素的最大公因子。这在任何欧几里得整环中都成立。两个元素的最大公因子的概念（若存在）可以在一般环中精确表述。

\begin{definition}
设 \(R\) 是交换环，\(a, b \in R\)（\(b \neq 0\)）。

（1）若存在元素 \(x \in R\) 使 \(a = bx\)，则称 \(a\) 是 \(b\) 的\textbf{倍数}。此时称 \(b\) \textbf{整除} \(a\)，或称 \(b\) 是 \(a\) 的\textbf{因子}，记作 \(b \mid a\).

（2）\(a\) 与 \(b\) 的\textbf{最大公因子}是非零元素 \(d\)，满足

（i）\(d \mid a\) 且 \(d \mid b\)，且

（ii）若 \(d' \mid a\) 且 \(d' \mid b\) 则 \(d' \mid d\).

\(a\) 与 \(b\) 的最大公因子记作 \(\gcd(a, b)\)，或（滥用记号）简单地记作 \((a, b)\).
\end{definition}

注意在环 \(R\) 中 \(b \mid a\) 当且仅当 \(a \in (b)\) 当且仅当 \((a) \subseteq (b)\). 特别地，若 \(d\) 是 \(a\) 和 \(b\) 的任意公因子，则 \((d)\) 必同时含 \(a\) 和 \(b\)，因而必含由 \(a\) 和 \(b\) 生成的理想。因此 \(a\) 与 \(b\) 的最大公因子的定义性质（i）和（ii）用理想的语言表述分别变成：
\begin{center}
若 \(I\) 是由 \(a\) 和 \(b\) 生成的 \(R\) 的理想，则 \(d\) 是 \(a\) 与 \(b\) 的最大公因子当且仅当\\
（i）\(I\) 包含在主理想 \((d)\) 中，且\\
（ii）若 \((d')\) 是任何包含 \(I\) 的主理想则 \((d) \subseteq (d')\).
\end{center}
因此 \(a\) 与 \(b\) 的最大公因子（若存在）是包含 \(a\) 和 \(b\) 的唯一的极小主理想的生成元。存在没有最大公因子的环。

这个讨论立即给出最大公因子存在的一个充分条件。

\begin{proposition}\label{pro:8.2}
若 \(a\) 和 \(b\) 是交换环 \(R\) 中的非零元素，使得由 \(a\) 和 \(b\) 生成的理想是主理想 \((d)\)，则 \(d\) 是 \(a\) 与 \(b\) 的最大公因子。
\end{proposition}

这解释了为什么记号 \((a, b)\) 常同时用来表示由 \(a\) 和 \(b\) 生成的理想以及 \(a\) 与 \(b\) 的最大公因子。每个由两个元素生成的理想 \((a, b)\) 都是主理想的整环称为 \textbf{Bézout 整环}。本节及后续各节的习题将探索这些环，并表明存在含有非主（必然无限生成）理想的 Bézout 整环。

注意命题 \ref{pro:8.2} 中的条件不是必要条件。例如在环 \(R = \mathbb{Z}[x]\) 中，元素 2 和 \(x\) 生成一个极大的非主理想（参见 7.4 节的例子）。因此 \(R = (1)\) 是同时包含 2 和 \(x\) 的唯一主理想，故 1 是 2 和 \(x\) 的最大公因子。我们将在 8.3 节看到这类其他例子。

在回到欧几里得整环之前，我们考察最大公因子的唯一性。

\begin{proposition}\label{pro:8.3}
设 \(R\) 是整环。若 \(R\) 的两个元素 \(d\) 和 \(d'\) 生成同一个主理想，即 \((d) = (d')\)，则对 \(R\) 的某个单位 \(u\) 有 \(d' = ud\). 特别地，若 \(d\) 和 \(d'\) 都是 \(a\) 与 \(b\) 的最大公因子，则对某个单位 \(u\) 有 \(d' = ud\).
\end{proposition}

\begin{proof}
若 \(d\) 或 \(d'\) 为零则显然，故可设 \(d, d'\) 非零。由于 \(d \in (d')\)，存在 \(x \in R\) 使 \(d = xd'\). 由于 \(d' \in (d)\)，存在 \(y \in R\) 使 \(d' = yd\). 故 \(d = xyd\)，从而 \(d(1-xy) = 0\). 由于 \(d \neq 0\)，\(xy = 1\)，即 \(x\) 和 \(y\) 都是单位。这证明了第一个论断。第二个论断由第一个推出，因为 \(a\) 与 \(b\) 的任意两个最大公因子生成同一个主理想（它们互相整除）。
\end{proof}

欧几里得整环最重要的性质之一是最大公因子总存在且可以算法地计算。

\begin{theorem}\label{thm:8.4}
设 \(R\) 是欧几里得整环，\(a, b\) 是 \(R\) 的非零元素。设 \(d = r_n\) 是本章开头描述的 \(a\) 与 \(b\) 的欧几里得算法中的最后一个非零余数。则

（1）\(d\) 是 \(a\) 与 \(b\) 的最大公因子，且

（2）主理想 \((d)\) 是由 \(a\) 和 \(b\) 生成的理想。特别地，\(d\) 可以写成 \(a\) 和 \(b\) 的 \(R\)-线性组合，即存在 \(R\) 中的元素 \(x, y\) 使
\[
d = ax + by.
\]
\end{theorem}

\begin{proof}
由命题 \ref{pro:8.1}，由 \(a\) 和 \(b\) 生成的理想是主的，故 \(a, b\) 确有最大公因子，即生成（主）理想 \((a, b)\) 的任意元素。因此只要我们证明 \(d = r_n\) 生成这个理想，定理的两部分都将推出，即只需证明

（i）\(d \mid a\) 且 \(d \mid b\)（故 \((a, b) \subseteq (d)\)）；

（ii）\(d\) 是 \(a\) 和 \(b\) 的 \(R\)-线性组合（故 \((d) \subseteq (a, b)\)）。

为证明 \(d\) 整除 \(a\) 和 \(b\)，只需跟踪欧几里得算法中的整除性。从第 \((n+1)\) 个等式 \(r_{n-1} = q_{n+1}r_n\) 出发，我们看到 \(r_n \mid r_{n-1}\). 显然 \(r_n \mid r_n\). 由归纳（从指标 \(n\) 向下到指标 0）假设 \(r_n\) 整除 \(r_{k+1}\) 和 \(r_k\). 由第 \((k+1)\) 个等式 \(r_{k-1} = q_{k+1}r_k + r_{k+1}\)，由于 \(r_n\) 整除右端两项，我们看到 \(r_n\) 也整除 \(r_{k-1}\). 由欧几里得算法中的第 1 个等式得到 \(r_n\) 整除 \(b\)，然后由第 0 个等式得到 \(r_n\) 整除 \(a\). 故（i）成立。

为证明 \(r_n\) 在由 \(a\) 和 \(b\) 生成的理想 \((a, b)\) 中，类似地对从等式（0）到等式（\(n\)）作归纳。由等式（0）得 \(r_0 \in (a, b)\)，由等式（1）得 \(r_1 = b - q_1r_0 \in (b, r_0) \subseteq (a, b)\). 由归纳假设 \(r_{k-1}, r_k \in (a, b)\). 则由第 \((k+1)\) 个等式
\[
r_{k+1} = r_{k-1} - q_{k+1}r_k \in (r_{k-1}, r_k) \subseteq (a, b).
\]
这个归纳表明 \(r_n \in (a, b)\)，完成证明。
\end{proof}

上面许多材料对整数 \(\mathbb{Z}\) 的情形可能从初等算术中已经熟悉，除了翻译成理想的语言。例如若 \(a = 2210\)、\(b = 1131\)，则 \(\mathbb{Z}\) 中同时包含 \(a\) 和 \(b\) 的最小理想（由 \(a\) 和 \(b\) 生成的理想）是 \(13\mathbb{Z}\)，因为 13 是 2210 与 1131 的最大公因子。这个事实由欧几里得算法快速得出：
\[
\begin{aligned}
2210 &= 1 \cdot 1131 + 1079\\
1131 &= 1 \cdot 1079 + 52\\
1079 &= 20 \cdot 52 + 39\\
52 &= 1 \cdot 39 + 13\\
39 &= 3 \cdot 13
\end{aligned}
\]
故 \(13 = (2210, 1131)\) 是最后一个非零余数。用定理 \ref{thm:8.4} 的方法我们也可以把 13 写成 2210 和 1131 的线性组合，方法是先由上面倒数第二个等式解出 \(13 = 52 - 1\cdot 39\)，然后用前面的等式解出 39 和 52 等，最后把 13 完全用 2210 和 1131 表示。此时答案是
\[
13 = (-22)\cdot 2210 + 43\cdot 1131.
\]

整数 \(\mathbb{Z}\) 中的欧几里得算法极快。有一个定理说确定两个整数 \(a\) 和 \(b\) 的最大公因子所需的步数至多是两个数中较小者位数的 5 倍。换句话说，这个算法关于整数大小的对数来说是线性的。为体会这里隐含的速度，注意对上面的例子我们至多预期 \(5\cdot 4 = 20\) 次除法（该例所需的远少于此）。如果我们从 \(10^{100}\) 量级的整数（按物理标准是大数）出发，我们至多只预期 500 次除法。

\((a, b) = ax + by\) 中的整数 \(x, y\) 没有唯一性陈述。事实上 \(x' = x + b\)、\(y' = y - a\) 满足 \((a, b) = ax' + by'\). 这基本上就是唯一的可能性——可以证明若 \(x_0, y_0\) 是方程 \(ax + by = N\) 的解，则这个方程的任何其他解 \(x, y\) 都形如
\[
x = x_0 + m\frac{b}{(a, b)}, \qquad y = y_0 - m\frac{a}{(a, b)}
\]
（\(m\) 为某个整数（正或负））。

后一个定理（其证明在习题中给出概要）在已知方程至少有一个解的前提下给出一阶丢番图方程 \(ax + by = N\) 的完整解。但方程 \(ax + by = N\) 只是 \(N\) 是由 \(a\) 和 \(b\) 生成的理想的元素的另一种说法。由于我们知道这个理想就是 \((d)\)（由 \(a\) 与 \(b\) 的最大公因子 \(d\) 生成的主理想），这与说 \(N \in (d)\) 即 \(N\) 被 \(d\) 整除是同一回事。因此方程 \(ax + by = N\) 在整数 \(x, y\) 中可解当且仅当 \(N\) 被 \(a\) 与 \(b\) 的最大公因子整除（此时上面引用的结果给出这个方程的完整解集）。

本节最后给出另一个有时可用来证明给定整环不是欧几里得整环的判别准则\footnote{这里和下一节部分材料遵循 J. C. Wilson 在 A principal ideal ring that is not a Euclidean ring, Math. Mag., 46(1973), pp. 34--38 中的论述（基于 Th. Motzkin 的思想），并使用了 Kenneth S. Williams 在 Note on non-Euclidean Principal Ideal Domains, Math. Mag., 48(1975), pp. 176--177 中的简化。}。对任意整环，用 \(R^* = R^\times \cup \{0\}\) 表示 \(R\) 的单位连同 0 的集合。元素 \(u \in R - R^*\) 称为\textbf{通用旁因子}，若对每个 \(x \in R\) 存在某个 \(z \in R^*\) 使得 \(u\) 在 \(R\) 中整除 \(x - z\)，即存在一种对 \(u\) 的“除法算法”：每个 \(x\) 都可以写成 \(x = qu + z\)，其中 \(z\) 是零或单位。通用旁因子的存在是欧几里得条件的一个弱化：

\begin{proposition}\label{pro:8.5}
设 \(R\) 是域以外的整环。若 \(R\) 是欧几里得整环，则 \(R\) 中存在通用旁因子。
\end{proposition}

\begin{proof}
假设 \(R\) 对某个范数 \(N\) 欧几里得，设 \(u\) 是 \(R - R^*\)（由于 \(R\) 不是域，这个集合非空）中有最小范数的元素。对任意 \(x \in R\)，写 \(x = qu + r\)，其中 \(r = 0\) 或 \(N(r) < N(u)\). 两种情形下 \(u\) 的极小性都蕴含 \(r \in R^*\). 故 \(u\) 是 \(R\) 中的通用旁因子。
\end{proof}

\begin{example}
我们可以用命题 \ref{pro:8.5} 证明二次整数环 \(R = \mathbb{Z}[(1+\sqrt{-19})/2]\) 对任何范数都不是欧几里得整环，方法是证明 \(R\) 不含通用旁因子（我们将在下一节看到 \(R\) 中所有理想都是主的，故命题 \ref{pro:8.1} 后的例子中的技巧不适用于这个环）。我们已经确定 \(\pm 1\) 是 \(R\) 中仅有的单位，故 \(R^* = \{0, \pm 1\}\). 假设 \(u \in R\) 是通用旁因子，并用 \(N(a + b(1+\sqrt{-19})/2) = a^2 + ab + 5b^2\) 表示 \(R\) 上的域范数（如 7.1 节）。注意若 \(a, b \in \mathbb{Z}\) 且 \(b \neq 0\)，则 \(a^2+ab+5b^2 = (a+b/2)^2 + \tfrac{19}{4}b^2 \geq 5\)，故 \(N\) 在 \(R\) 上的最小非零值是 1（对单位 \(\pm 1\)）和 4（对 \(\pm 2\)）。在通用旁因子的定义中取 \(x = 2\)，可知 \(u\) 必须在 \(R\) 中整除 \(2-0\) 或 \(2\pm 1\) 之一，即 \(u\) 是 2 或 3 在 \(R\) 中的非单位因子。若 \(2 = \alpha\beta\) 则 \(4 = N(\alpha)N(\beta)\)，由上面的评注可知 \(\alpha\) 或 \(\beta\) 的范数为 1，即等于 \(\pm 1\). 故 2 在 \(R\) 中仅有的因子是 \(\{\pm 1, \pm 2\}\). 类似地，3 在 \(R\) 中仅有的因子是 \(\{\pm 1, \pm 3\}\)，故 \(u\) 仅有的可能值是 \(\pm 2\) 或 \(\pm 3\). 取 \(x = (1+\sqrt{-19})/2\)，容易验证 \(x, x\pm 1\) 都不被 \(R\) 中的 \(\pm 2\) 或 \(\pm 3\) 整除，故这些都不是通用旁因子。
\end{example}

\subsection*{习题}

\begin{enumerate}
\item 对下列五对整数 \(a, b\)，求出它们的最大公因子 \(d\) 并把 \(d\) 写成 \(a\) 和 \(b\) 的线性组合 \(ax+by\).

（a）\(a = 20, b = 13\).（b）\(a = 69, b = 372\).（c）\(a = 11391, b = 5673\).（d）\(a = 507885, b = 60808\).（e）\(a = 91442056588823, b = 779086434385541\)（对这些整数欧几里得算法只需 7 步）。

\item 对下列每对整数 \(a, n\)，证明 \(a\) 与 \(n\) 互素并求出 \(a\) 模 \(n\) 的逆元（参见预备知识第 3 节）。

（a）\(a = 13, n = 20\).（b）\(a = 69, n = 89\).（c）\(a = 1891, n = 3797\).（d）\(a = 6003722857, n = 77695236973\)（对这些整数欧几里得算法只需 3 步）。

\item 设 \(R\) 是欧几里得整环。设 \(m\) 是 \(R\) 的非零元素的范数集合中的最小整数。证明 \(R\) 中范数为 \(m\) 的每个非零元素都是单位。由此推出范数为零的非零元素（若存在）是单位。

\item 设 \(R\) 是欧几里得整环。

（a）证明：若 \((a, b) = 1\) 且 \(a\) 整除 \(bc\)，则 \(a\) 整除 \(c\). 更一般地，证明：若 \(a\) 整除 \(bc\)（\(a, b\) 非零），则 \(\frac{a}{(a, b)}\) 整除 \(c\).

（b）考虑丢番图方程 \(ax + by = N\)，其中 \(a, b, N\) 是整数且 \(a, b\) 非零。假设 \(x_0, y_0\) 是解：\(ax_0 + by_0 = N\). 证明这个方程的完整解集由
\[
x = x_0 + m\frac{b}{(a, b)}, \qquad y = y_0 - m\frac{a}{(a, b)}
\]
给出（\(m\) 取遍整数）。〔若 \(x, y\) 是 \(ax + by = N\) 的解，证明 \(a(x-x_0) = b(y_0-y)\) 并用（a）。〕

\item 求下列方程的所有整数解：

（a）\(2x + 4y = 5\)；（b）\(17x + 29y = 31\)；（c）\(85x + 145y = 505\).

\item （\textbf{邮票问题}）设 \(a, b\) 是两个互素的正整数。证明每个充分大的正整数 \(N\) 都可以写成 \(a\) 和 \(b\) 的线性组合 \(ax + by\)，其中 \(x, y\) 都非负，即存在整数 \(N_0\) 使对所有 \(N \geq N_0\) 方程 \(ax + by = N\) 都可以用非负整数 \(x, y\) 求解。证明事实上整数 \(ab - a - b\) 不能写成 \(a\) 和 \(b\) 的正线性组合，但每个大于 \(ab-a-b\) 的整数都是 \(a\) 和 \(b\) 的正线性组合（故每个大于 \(ab-a-b\) 的“邮资”都可以只用面值 \(a, b\) 的邮票得到）。

\item 用欧几里得算法求 \(\mathbb{Z}[i]\) 中理想 \((85, 1+13i)\) 的生成元，即 85 与 \(1+13i\) 的最大公因子。对理想 \((47-13i, 53+56i)\) 做同样的事。

\noindent 已知（但不太容易证明）\(D = -1, -2, -3, -7, -11, -19, -43, -67, -163\) 是使 \(\mathcal{O}\) 中每个理想都是主理想（即下一节术语中的 \(\mathcal{O}\) 是 P.I.D.）的仅有的 \(D\) 的负值。下一习题的结果精确确定哪些 \(D < 0\) 的二次整数环是欧几里得的。

\item 设 \(F = \mathbb{Q}(\sqrt{D})\) 是二次域，相关的二次整数环为 \(\mathcal{O}\)，域范数 \(N\) 如 7.1 节。

（a）假设 \(D\) 是 \(-1, -2, -3, -7\) 或 \(-11\). 证明 \(\mathcal{O}\) 关于 \(N\) 是欧几里得整环。〔修改正文中对 \(\mathbb{Z}[i]\)（\(D = -1\)）的证明。对 \(D = -3, -7, -11\)，证明 \(F\) 的每个元素与 \(\mathcal{O}\) 中某个元素之差是一个范数至多为 \((1+|D|)^2/(16|D|)\) 的元素，而这对这些 \(D\) 值都小于 1. 在 \(\mathbb{C}\) 中画出 \(\mathcal{O}\) 的点可能有帮助。〕

（b）假设 \(D = -43, -67\) 或 \(-163\). 证明 \(\mathcal{O}\) 对任何范数都不是欧几里得整环。〔应用正文中对 \(D = -19\) 的同样证明。〕

\item 证明二次域 \(\mathbb{Q}(\sqrt{2})\) 中的整数环 \(\mathcal{O}\) 关于由 7.1 节域范数 \(N\) 的绝对值给出的范数是欧几里得整环。

\item 证明对 \(\mathbb{Z}[i]\) 的任意非零理想 \(I\)，商环 \(\mathbb{Z}[i]/I\) 有限。〔用 \(I = (a)\)（\(a\) 非零）这一事实，然后用这个欧几里得整环中的除法算法看出 \(I\) 的每个陪集都由范数小于 \(N(a)\) 的元素代表。〕

\item 设 \(R\) 是带 1 的交换环，\(a, b\) 是 \(R\) 的非零元素。\(a\) 与 \(b\) 的\textbf{最小公倍}是 \(R\) 的元素 \(e\)，满足

（i）\(a \mid e\) 且 \(b \mid e\)，且

（ii）若 \(a \mid e'\) 且 \(b \mid e'\) 则 \(e \mid e'\).

（a）证明 \(a\) 与 \(b\) 的最小公倍（若存在）是含于 \((a) \cap (b)\) 中的唯一的极大主理想的生成元。

（b）由此推出欧几里得整环中任意两个非零元素都有最小公倍，且它在乘以单位的意义下唯一。

（c）证明欧几里得整环中 \(a\) 与 \(b\) 的最小公倍是 \(\dfrac{ab}{(a, b)}\)，其中 \((a, b)\) 是 \(a\) 与 \(b\) 的最大公因子。

\item （\textbf{公钥密码}）设 \(N\) 是正整数。设 \(M\) 是与 \(N\) 互素的整数，\(d\) 是与 \(\varphi(N)\) 互素的整数，其中 \(\varphi\) 表示欧拉 \(\varphi\)-函数。证明若 \(M_1 \equiv M^d \pmod N\)，则 \(M \equiv M_1^{d'} \pmod N\)，其中 \(d'\) 是 \(d\) 模 \(\varphi(N)\) 的逆元：\(dd' \equiv 1 \pmod{\varphi(N)}\).

\noindent\textbf{注：}这个结果是标准公钥密码的基础。假设 \(N = pq\) 是两个不同大素数之积（例如每个约 100 位）。若 \(M\) 是消息，则 \(M_1 \equiv M^d \pmod N\) 是 \(M\) 的扰乱（编码）版本，可以通过计算 \(M_1^{d'} \pmod N\) 来还原（解码）（这些幂即使对大的 \(M, N\) 也可以通过反复平方相当容易地计算）。\(N\) 和 \(d\)（但不是 \(p, q\)）被公开（故名），然后任何有消息的人都可以发送他们编码的消息 \(M^d \pmod N\). 解码消息似乎需要确定 \(d'\)，而这需要确定 \(\varphi(N) = \varphi(pq) = (p-1)(q-1)\) 的值（然而至今没有人证明不存在其他解码方法）。这个方法作为密码的成功依赖于确定 \(N\) 的素数分解的必要性，而对此没有足够高效的算法（例如最朴素的方法——检查到 \(\sqrt{N}\) 的所有因子——在这里需要约 \(10^{100}\) 次计算，即使每秒 100 亿次计算也需极其漫长的时间；当然，人们总可以增大 \(p, q\) 的规模）。
\end{enumerate}

\section{主理想整环（P.I.D.）}

\begin{definition}
\textbf{主理想整环}（P.I.D.）是每个理想都是主理想的整环。
\end{definition}

命题 \ref{pro:8.1} 证明了每个欧几里得整环都是主理想整环，故关于主理想整环的每个结果自动对欧几里得整环成立。

\begin{example}
（1）如命题 \ref{pro:8.1} 后所述，整数 \(\mathbb{Z}\) 是 P.I.D. 我们在 7.4 节看到多项式环 \(\mathbb{Z}[x]\) 含非主理想，故不是 P.I.D.

（2）命题 \ref{pro:8.1} 后的例 2 表明二次整数环 \(\mathbb{Z}[\sqrt{-5}]\) 不是 P.I.D.，事实上理想 \((3, 1+\sqrt{-5})\) 是非主理想。两个非主理想 \(I\) 和 \(J\) 的积有可能是主理想，例如理想 \((3, 1+\sqrt{-5})\) 和 \((3, 1-\sqrt{-5})\) 都是非主的，但它们的积是由 3 生成的主理想，即 \((3, 1+\sqrt{-5})(3, 1-\sqrt{-5}) = (3)\)（参见习题 5 以及下面命题 \ref{pro:8.9} 前的例子）。

并非每个主理想整环都是欧几里得整环。我们将在下面证明二次整数环 \(\mathbb{Z}[(1+\sqrt{-19})/2]\)（上一节已证明它不是欧几里得整环）却是 P.I.D.

从理想论的观点看，主理想整环是域（其理想只有平凡的 (0) 和 (1)）之外自然要研究的一类环。欧几里得整环享有的许多性质主理想整环也满足。然而欧几里得整环相对于主理想整环的一个重要优点是：虽然两种情形下最大公因子都存在，但在欧几里得整环中人们有计算它们的算法。因此（特别是我们将在第 12 章看到的）依赖最大公因子存在性的结果常常可以在更大的主理想整环类中证明，尽管例子（即这些结果的具体应用）的计算在（若有）欧几里得算法时更有效。
\end{example}

我们汇总上一节证明的关于最大公因子的一些事实。

\begin{proposition}\label{pro:8.6}
设 \(R\) 是主理想整环，\(a, b\) 是 \(R\) 的非零元素。设 \(d\) 是由 \(a\) 和 \(b\) 生成的主理想的生成元。则

（1）\(d\) 是 \(a\) 与 \(b\) 的最大公因子；

（2）\(d\) 可以写成 \(a\) 和 \(b\) 的 \(R\)-线性组合，即存在 \(R\) 中的元素 \(x, y\) 使
\[
d = ax + by;
\]

（3）\(d\) 在乘以 \(R\) 的单位的意义下唯一。
\end{proposition}

\begin{proof}
这只是命题 \ref{pro:8.2} 和 \ref{pro:8.3}.
\end{proof}

回忆极大理想总是素理想，但反过来一般不对。然而我们在 7.4 节观察到 \(\mathbb{Z}\) 的每个非零素理想都是极大理想。这个有用的事实对任意主理想整环成立，如下一命题所示。

\begin{proposition}\label{pro:8.7}
主理想整环中每个非零素理想都是极大理想。
\end{proposition}

\begin{proof}
设 \((p)\) 是主理想整环 \(R\) 中的非零素理想，\(I = (m)\) 是任意包含 \((p)\) 的理想。我们必须证明 \(I = (p)\) 或 \(I = R\). 现在 \(p \in (m)\)，故对某个 \(r \in R\) 有 \(p = rm\). 由于 \((p)\) 是素理想且 \(rm \in (p)\)，\(r\) 或 \(m\) 必在 \((p)\) 中。若 \(m \in (p)\) 则 \((p) = (m) = I\). 另一方面若 \(r \in (p)\)，写 \(r = ps\). 此时 \(p = rm = psm\)，故 \(sm = 1\)（回忆 \(R\) 是整环），\(m\) 是单位，故 \(I = R\).
\end{proof}

如我们已经提到的，若 \(F\) 是域，则多项式环 \(F[x]\) 是欧几里得整环，因而也是主理想整环（这将在下一章证明）。其逆也成立。直观上，若 \(I\) 是 \(R\) 中的理想（如 \(\mathbb{Z}\) 中的理想 \((2)\)），则 \(R[x]\) 中的理想 \((I, x)\)（如 \(\mathbb{Z}[x]\) 中的理想 \((2, x)\)）比 \(I\) 多需要一个生成元，故一般不是主的。

\begin{corollary}\label{cor:8.8}
若 \(R\) 是使多项式环 \(R[x]\) 为主理想整环（或欧几里得整环）的任意交换环，则 \(R\) 必是域。
\end{corollary}

\begin{proof}
假设 \(R[x]\) 是主理想整环。由于 \(R\) 是 \(R[x]\) 的子环，\(R\) 必是整环（回忆 \(R[x]\) 有单位元当且仅当 \(R\) 有）。理想 \((x)\) 是 \(R[x]\) 中的非零素理想，因为 \(R[x]/(x)\) 同构于整环 \(R\). 由命题 \ref{pro:8.7}，\((x)\) 是极大理想，故由 7.4 节命题 12，商 \(R\) 是域。
\end{proof}

本节最后一个结果将用于证明并非每个 P.I.D. 都是欧几里得整环，并把主理想性质与欧几里得条件的另一个弱化联系起来。

\begin{definition}
称 \(N\) 是 \textbf{Dedekind--Hasse 范数}，若 \(N\) 是正范数且对每个非零 \(a, b \in R\)，或者 \(a\) 是理想 \((b)\) 的元素，或者理想 \((a, b)\) 中存在范数严格小于 \(b\) 的范数的非零元素（即或者在 \(R\) 中 \(b\) 整除 \(a\)，或者存在 \(s, t \in R\) 使 \(0 < N(sa - tb) < N(b)\)）。
\end{definition}

注意若总能取 \(s = 1\) 满足 Dedekind--Hasse 条件，则 \(R\) 关于正范数 \(N\) 欧几里得，故这确实是欧几里得条件的一个弱化。

\begin{proposition}\label{pro:8.9}
整环 \(R\) 是 P.I.D. 当且仅当 \(R\) 有 Dedekind--Hasse 范数。
\end{proposition}

\begin{proof}
设 \(I\) 是 \(R\) 的任意非零理想，\(b\) 是 \(I\) 中使 \(N(b)\) 最小的非零元素。假设 \(a\) 是 \(I\) 中的任意非零元素，故理想 \((a, b)\) 包含在 \(I\) 中。则 \(N\) 上的 Dedekind--Hasse 条件以及 \(b\) 的极小性蕴含 \(a \in (b)\)，故 \(I = (b)\) 是主的。逆命题将在下一节证明（推论 \ref{cor:8.16}）。
\end{proof}

\begin{example}
设 \(R = \mathbb{Z}[(1+\sqrt{-19})/2]\) 是上一节末考虑的二次整数环。我们证明定义在 \(R\) 上的正域范数 \(N(a + b(1+\sqrt{-19})/2) = a^2+ab+5b^2\) 是 Dedekind--Hasse 范数，由命题 \ref{pro:8.9} 以及上一节的结果这将证明 \(R\) 是 P.I.D. 但不是欧几里得整环。

假设 \(\alpha, \beta\) 是 \(R\) 的非零元素且 \(\alpha/\beta \notin R\). 我们必须证明存在 \(s, t \in R\) 使 \(0 < N(s\alpha - t\beta) < N(\beta)\)，由域范数的乘性这等价于
\[
0 < N\left(s - t\frac{\alpha}{\beta}\right) < 1. \tag{$*$}
\]
把 \(\frac{\alpha}{\beta} = \frac{a+b\sqrt{-19}}{c}\) 写成具有整数 \(a, b, c\)、无公因子且 \(c > 1\) 的形式（因为假设 \(\beta\) 不整除 \(\alpha\)）。由于 \(a, b, c\) 无公因子，存在整数 \(x, y, z\) 使 \(ax+by+cz = 1\). 对某些商 \(q\) 和余数 \(r\)（\(|r| \leq c/2\)）写 \(ay - 19bx = cq + r\)，令 \(s = y + x\sqrt{-19}\)、\(t = q - z\sqrt{-19}\). 则一个快速计算表明
\[
0 < N\left(s - t\frac{\alpha}{\beta}\right) = \frac{(ay-19bx-cq)^2 + 19(ax+by+cz)^2}{4c^2} \leq \frac{1}{4} + \frac{19}{4c^2},
\]
故只要 \(c \geq 5\)，这个 \(s, t\) 就满足 \((*)\).

假设 \(c = 2\). 则 \(a, b\) 中一个是偶数、另一个是奇数（否则 \(\frac{\alpha}{\beta} \in R\)），此时容易验证 \(s = 1\) 和 \(t = \frac{(a-1)+b\sqrt{-19}}{2}\) 是满足 \((*)\) 的 \(R\) 中元素。

假设 \(c = 3\). 整数 \(a^2+19b^2\) 不被 3 整除（模 3 这是 \(a^2+b^2\)，容易看出它模 3 为 0 当且仅当 \(a, b\) 都模 3 为 0；但那样 \(a, b, c\) 有公因子）。写 \(a^2+19b^2 = 3q + r\)，其中 \(r = 1\) 或 2. 再次容易验证 \(s = a - b\sqrt{-19}\)、\(t = q\) 是满足 \((*)\) 的 \(R\) 中元素。

最后假设 \(c = 4\)，故 \(a, b\) 不都是偶数。若 \(a, b\) 中一个是偶数、另一个是奇数，则 \(a^2+19b^2\) 是奇数，故我们可以对某些 \(q, r \in \mathbb{Z}\)（\(0 < r < 4\)）写 \(a^2+19b^2 = 4q+r\). 则 \(s = a-b\sqrt{-19}\)、\(t = q\) 满足 \((*)\). 若 \(a, b\) 都是奇数，则 \(a^2+19b^2 \equiv 1+3 \equiv 4 \pmod 8\)，故我们可以对某个 \(q \in \mathbb{Z}\) 写 \(a^2+19b^2 = 8q+4\). 则 \(s = \frac{a-b\sqrt{-19}}{2}\)、\(t = q\) 是满足 \((*)\) 的 \(R\) 中元素。
\end{example}

\subsection*{习题}

\begin{enumerate}
\item 证明主理想整环中两个理想 \((a)\) 和 \((b)\) 互余（参见 7.6 节）当且仅当 \(a\) 与 \(b\) 的最大公因子是 1（此时称 \(a\) 与 \(b\) \textbf{互素}）。

\item 证明 P.I.D. 中任意两个非零元素都有最小公倍（参见第 1 节习题 11）。

\item 证明 P.I.D. 对素理想的商仍是 P.I.D.

\item 设 \(R\) 是整环。证明若下列两个条件成立则 \(R\) 是主理想整环：

（i）\(R\) 中任意两个非零元素 \(a, b\) 都有最大公因子，且它可以写成 \(ra + sb\) 的形式（\(r, s \in R\)）；

（ii）若 \(a_1, a_2, a_3, \ldots\) 是 \(R\) 的非零元素且对所有 \(i\) 有 \(a_{i+1} \mid a_i\)，则存在正整数 \(N\) 使得对所有 \(n \geq N\)，\(a_n\) 是 \(a_N\) 的单位倍。

\item 设 \(R\) 是二次整数环 \(\mathbb{Z}[\sqrt{-5}]\). 定义理想 \(I_1 = (2, 1+\sqrt{-5})\)、\(I_2 = (3, 2+\sqrt{-5})\) 和 \(I_3 = (3, 2-\sqrt{-5})\).

（a）证明 \(I_1, I_2\) 和 \(I_3\) 是 \(R\) 中的非主理想。〔注意命题 \ref{pro:8.1} 后的例 2 已对 \(I_2\) 证明了这一点。〕

（b）通过证明 \(I_1^2\) 是由 2 生成的主理想（即 \(I_1^2 = (2)\)）来证明两个非主理想的积可以是主的。

（c）类似地证明 \(I_1I_2 = (1-\sqrt{-5})\) 和 \(I_1I_3 = (1+\sqrt{-5})\) 是主的。由此推出主理想 (6) 是 4 个理想的积：\((6) = I_1^2 I_2 I_3\).

\item 设 \(R\) 是整环，假设 \(R\) 中每个素理想都是主理想。这个习题证明 \(R\) 的每个理想都是主理想，即 \(R\) 是 P.I.D.

（a）假设 \(R\) 中非主理想的集合非空，证明这个集合按包含关系有极大元（由假设它不是素理想）。〔用佐恩引理。〕

（b）设 \(I\) 是在“非主”这一性质下极大的理想，\(a, b \in R\) 满足 \(ab \in I\) 但 \(a \notin I\)、\(b \notin I\). 设 \(I_a = (I, a)\) 是由 \(I\) 和 \(a\) 生成的理想，\(I_b = (I, b)\) 是由 \(I\) 和 \(b\) 生成的理想，定义 \(J = \{r \in R \mid ra \in I\}\). 证明 \(I_a = (\alpha)\) 和 \(J = (\beta)\) 是 \(R\) 中的主理想，且 \(I \subseteq I_a \subseteq J\)、\(I_aJ = (\alpha\beta) \subseteq I\).

（c）若 \(x \in I\)，证明对某个 \(s \in J\) 有 \(x = s\alpha\). 由此推出 \(I = I_a\) 是主理想，矛盾，并得出 \(R\) 是 P.I.D.

\item 每个由两个元素生成的理想都是主理想（即对每个 \(a, b \in R\)，\((a, b) = (d)\) 对某个 \(d \in R\)）的整环 \(R\) 称为 \textbf{Bézout 整环}。〔另见第 3 节习题 11。〕

（a）证明整环 \(R\) 是 Bézout 整环当且仅当 \(R\) 的每对元素 \(a, b\) 都有最大公因子 \(d \in R\)，且它可以写成 \(a\) 和 \(b\) 的 \(R\)-线性组合，即对某些 \(x, y \in R\) 有 \(d = ax+by\).

（b）证明 Bézout 整环的每个有限生成理想都是主理想。〔关于并非每个理想都是主理想的 Bézout 整环，参见 9.2 和 9.3 节的习题。〕

（c）设 \(F\) 是 Bézout 整环 \(R\) 的分式域。证明 \(F\) 的每个元素都可以写成 \(a/b\) 的形式，其中 \(a, b \in R\) 且 \(a, b\) 互素（参见习题 1）。

\item 证明：若 \(R\) 是主理想整环且 \(D\) 是 \(R\) 的闭于乘法的子集，则 \(D^{-1}R\) 也是 P.I.D.（参见 7.5 节）。
\end{enumerate}

\section{唯一分解整环（U.F.D.）}

对整数 \(\mathbb{Z}\)，还有另一种从初等算术中熟悉的、确定两个元素 \(a\) 和 \(b\) 的最大公因子的方法，即 \(a\) 和 \(b\) 的“分解为素数”的概念，由它可以容易地确定最大公因子。这也可以推广到一类更大的环，称为\textbf{唯一分解整环}（U.F.D.）——它们将很快定义。我们随后将证明
\begin{center}
每个主理想整环都是唯一分解整环，
\end{center}
故关于唯一分解整环的每个结果自动对欧几里得整环和主理想整环都成立。

我们首先引入一些术语。

\begin{definition}
设 \(R\) 是整环。

（1）假设 \(r \in R\) 非零且不是单位。则称 \(r\) 在 \(R\) 中\textbf{不可约}，若只要 \(r = ab\)（\(a, b \in R\)），\(a\) 或 \(b\) 中至少有一个必是 \(R\) 的单位。否则称 \(r\) 是\textbf{可约的}。

（2）非零元素 \(p \in R\) 称为 \(R\) 中的\textbf{素元}，若由 \(p\) 生成的理想 \((p)\) 是素理想。换句话说，非零元素 \(p\) 是素元，若它不是单位且只要对任意 \(a, b \in R\) 有 \(p \mid ab\)，就有 \(p \mid a\) 或 \(p \mid b\).

（3）\(R\) 的两个相差一个单位的元素 \(a\) 和 \(b\) 称为在 \(R\) 中\textbf{相伴}（即对 \(R\) 的某个单位 \(u\) 有 \(a = ub\)）。
\end{definition}

\begin{proposition}\label{pro:8.10}
在整环中素元总是不可约的。
\end{proposition}

\begin{proof}
假设 \((p)\) 是非零素理想且 \(p \mid ab\). 则 \(ab = p \in (p)\)，故由素理想的定义，\(a\) 或 \(b\)（设为 \(a\)）在 \((p)\) 中。于是对某个 \(r\) 有 \(a = pr\). 这蕴含 \(p = ab = prb\)，故 \(rb = 1\)，\(b\) 是单位。这表明 \(p\) 不可约。
\end{proof}

一般地，不可约元素不必是素元。例如考虑二次整数环 \(R = \mathbb{Z}[\sqrt{-5}]\) 中的元素 3. 8.1 节的计算表明 3 在 \(R\) 中不可约，但 3 不是素元，因为 \((2+\sqrt{-5})(2-\sqrt{-5}) = 3^2\) 被 3 整除，但 \(2+\sqrt{-5}\) 和 \(2-\sqrt{-5}\) 在 \(R\) 中都不被 3 整除。

然而若 \(R\) 是主理想整环，则素元与不可约元素的概念相同。特别地，这些概念在 \(\mathbb{Z}\) 和 \(F[x]\)（\(F\) 是域）中重合。

\begin{proposition}\label{pro:8.11}
在主理想整环中，非零元素是素元当且仅当它不可约。
\end{proposition}

\begin{proof}
上面我们已经证明素元蕴含不可约。我们必须反过来证明：若 \(p\) 不可约，则 \(p\) 是素元，即理想 \((p)\) 是素理想。若 \(M\) 是任意包含 \((p)\) 的理想，则由假设 \(M = (m)\) 是主理想。由于 \(p \in (m)\)，对某个 \(r\) 有 \(p = rm\). 但 \(p\) 不可约，故按定义 \(r\) 或 \(m\) 是单位。这分别意味着 \((p) = (m)\) 或 \((m) = (1)\). 因此包含 \((p)\) 的仅有理想是 \((p)\) 或 \((1)\)，即 \((p)\) 是极大理想。由于极大理想是素理想，证明完成。
\end{proof}

\begin{example}
命题 \ref{pro:8.11} 给出另一个证明二次整数环 \(\mathbb{Z}[\sqrt{-5}]\) 不是 P.I.D. 的方法，因为 3 在这个环中不可约但不是素元。
\end{example}

整数 \(\mathbb{Z}\) 中的不可约元素是初等算术中熟悉的素数（及其相反数），且两个整数 \(a\) 和 \(b\) 相伴当且仅当 \(a = \pm b\).

在整数 \(\mathbb{Z}\) 中，每个整数 \(n\) 都可以写成素数的乘积（不必互异），如下。若 \(n\) 本身不是素数，则按定义可以写成 \(n = n_1n_2\)，其中 \(n_1\) 和 \(n_2\) 是两个都不是单位的整数，即都不是 \(\pm 1\). \(n_1\) 和 \(n_2\) 的绝对值都必小于 \(n\) 本身。若它们都是素数，我们已经把 \(n\) 写成了素数之积。若 \(n_1\) 或 \(n_2\) 中有一个不是素数，则它又可以分解为两个（更小的）整数。由于整数的绝对值不能无限减小，我们必在某个时刻只剩下素数因子，于是我们把 \(n\) 写成了素数之积。

例如若 \(n = 2210\)，上面的算法如下进行：\(n\) 本身不是素数，因为我们可以写 \(n = 2 \cdot 1105\). 整数 2 是素数，但 1105 不是：\(1105 = 5 \cdot 221\). 整数 5 是素数，但 221 不是：\(221 = 13 \cdot 17\). 这里算法终止，因为 13 和 17 都是素数。这给出 2210 的素因子分解：\(2210 = 2 \cdot 5 \cdot 13 \cdot 17\). 类似地我们求出 \(1131 = 3 \cdot 13 \cdot 29\). 这些例子中每个素数只出现一次幂，但当然一般不必如此。

在环 \(\mathbb{Z}\) 中，不仅每个整数 \(n\) 都可以写成素数之积，事实上这个分解在如下意义下是唯一的：同一正整数 \(n\) 的任意两个素因子分解只相差正素因子书写的顺序。限制为正整数是为了避免把 15 的分解 \((3)(5)\) 和 \((-3)(-5)\) 视为本质不同。\(\mathbb{Z}\) 的这个唯一分解性质（我们将很快证明）对整数的算术极为有用。具有类似性质的一般环被赋予一个名称。

\begin{definition}
\textbf{唯一分解整环}（U.F.D.）是整环 \(R\)，其中每个非零且不是单位的元素 \(r \in R\) 具有下列两个性质：

（i）\(r\) 可以写成 \(R\) 的不可约元素 \(p_i\) 的有限乘积（不必互异）：\(r = p_1p_2\cdots p_n\)，且

（ii）（i）中的分解在相伴意义下唯一：即若 \(r = q_1q_2\cdots q_m\) 是 \(r\) 的另一个不可约分解，则 \(m = n\)，且因子的某种重新编号使 \(p_i\) 与 \(q_i\) 相伴（\(i = 1, 2, \ldots, n\)）。
\end{definition}

\begin{example}
（1）域 \(F\) 平凡地是唯一分解整环，因为每个非零元素都是单位，故没有需要验证性质（i）和（ii）的元素。

（2）如上面所指出的，我们将很快证明每个主理想整环都是唯一分解整环（特别地，\(\mathbb{Z}\) 和 \(F[x]\)（\(F\) 是域）都是唯一分解整环）。

（3）我们还将在下一章证明：只要 \(R\) 本身是唯一分解整环，多项式环 \(R[x]\) 就是唯一分解整环（这与主理想整环或欧几里得整环的性质形成对比，后两者不能从环 \(R\) 传递到多项式环 \(R[x]\)）。这个结果连同前一例子将表明 \(\mathbb{Z}[x]\) 是唯一分解整环。

（4）高斯整数的子环 \(R = \mathbb{Z}[2i] = \{a + 2bi \mid a, b \in \mathbb{Z}\}\)（其中 \(i^2 = -1\)）是整环但不是唯一分解整环（这类环在 7.1 节习题 23 中引入）。元素 2 和 \(2i\) 是 \(R\) 中不相伴的不可约元素（因为 \(i \notin R\)），且 \(4 = 2 \cdot 2 = (-2i)\cdot(2i)\) 是 \(R\) 中 4 的两个不同分解。也可以直接验证 \(2i\) 在 \(R\) 中不可约但不是素元（因为 \(R/(2i) \cong \mathbb{Z}/4\mathbb{Z}\)）。在更大的高斯整数环 \(\mathbb{Z}[i]\)（它是唯一分解整环）中，2 和 \(2i\) 相伴，因为 \(i\) 是这个更大环中的单位。我们将在 9.3 节末给出 \(\mathbb{Z}[2i]\) 不是唯一分解整环的稍有不同的证明（其中不必验证 2 和 \(2i\) 不可约）。

（5）二次整数环 \(\mathbb{Z}[\sqrt{-5}]\) 是另一个不是唯一分解整环的整环的例子，因为 \(6 = 2\cdot 3 = (1+\sqrt{-5})(1-\sqrt{-5})\) 给出 6 的两个不同不可约分解。\(\mathbb{Z}[\sqrt{-5}]\) 中的主理想 (6) 可以写成 4 个非主素理想之积：\((6) = P_2^2P_3P_3'\)，而 \(\mathbb{Z}[\sqrt{-5}]\) 中元素 6 的两个不同分解可以解释为把这个理想的积重新排列成主理想之积的两种方式：\(P_2^2 = (2)\) 与 \(P_3P_3' = (3)\) 之积，以及 \(P_2P_3 = (1+\sqrt{-5})\) 与 \(P_2P_3' = (1-\sqrt{-5})\) 之积（参见习题 8）。

虽然二次整数环 \(\mathcal{O}\) 的元素不必有唯一分解，但有一个定理（推论 16.16）说 \(\mathcal{O}\) 中每个理想都可以唯一地写成素理想之积。理想唯一分解为素理想之积对代数数域的整数环（二次整数环是其例子）一般成立，并引出第 16 章考虑的 Dedekind 整环的概念。正是数论中代数整数环的元素缺乏唯一分解，最初导致了理想的定义。这些环中分解为素理想的唯一性使理想的元素具有一种“理想的”（在“完美的”或“合乎需要的”意义上）行为，这是这些（现在已是基本的）代数对象术语选择的依据。
\end{example}

唯一分解整环中不可约元素的第一个性质是它们也是素元。人们可能以为我们可以从这个命题连同前面提到的定理（我们将很快证明：每个主理想整环都是唯一分解整环）推出命题 \ref{pro:8.11}，然而命题 \ref{pro:8.11} 将被用于后一定理的证明。

\begin{proposition}\label{pro:8.12}
在唯一分解整环中，非零元素是素元当且仅当它不可约。
\end{proposition}

\begin{proof}
设 \(R\) 是唯一分解整环。由命题 \ref{pro:8.10}，\(R\) 的素元不可约，故剩下要证明每个不可约元素都是素元。设 \(p\) 是 \(R\) 中的不可约元素，假设对某些 \(a, b \in R\) 有 \(p \mid ab\)；我们必须证明 \(p\) 整除 \(a\) 或 \(b\). 说 \(p\) 整除 \(ab\) 就是说对 \(R\) 中某个 \(c\) 有 \(ab = pc\). 把 \(a\) 和 \(b\) 写成不可约元素之积，由后一个等式以及 \(ab\) 的不可约分解的唯一性，我们看到不可约元素 \(p\) 必与出现在 \(a\) 的分解或 \(b\) 的分解中的某个不可约元素相伴。我们可以假设 \(p\) 与出现在 \(a\) 的分解中的一个不可约元素相伴（否则就交换 \(a, b\) 的角色）。换句话说，我们可以把 \(a\) 写成 \(a = (up)p_2\cdots p_n\)（\(u\) 是单位）的形式，因此 \(p\) 整除 \(a\)，因为 \(a = pd\)，其中 \(d = up_2\cdots p_n\)，完成证明。
\end{proof}

在唯一分解整环中，我们现在将不加区别地使用“素元”和“不可约元素”这两个术语，尽管我们通常称 \(\mathbb{Z}\) 中的“素数”和 \(F[x]\) 中的“不可约多项式”。

我们将用前一命题证明：在唯一分解整环中任意两个非零元素 \(a\) 和 \(b\) 都有最大公因子：

\begin{proposition}\label{pro:8.13}
设 \(a\) 和 \(b\) 是唯一分解整环 \(R\) 的两个非零元素，并假设
\[
a = up_1^{e_1}p_2^{e_2}\cdots p_n^{e_n} \quad \text{和} \quad b = vp_1^{f_1}p_2^{f_2}\cdots p_n^{f_n}
\]
是 \(a\) 和 \(b\) 的素因子分解，其中 \(u\) 和 \(v\) 是单位，素数 \(p_1, p_2, \ldots, p_n\) 互不相同，指数 \(e_i, f_i \geq 0\). 则元素
\[
d = p_1^{\min(e_1, f_1)}p_2^{\min(e_2, f_2)}\cdots p_n^{\min(e_n, f_n)}
\]
（若所有指数都为 0 则 \(d = 1\)）是 \(a\) 与 \(b\) 的最大公因子。
\end{proposition}

\begin{proof}
由于 \(d\) 中出现的每个素数的指数都不大于它在 \(a\) 和 \(b\) 的分解中出现的指数，\(d\) 同时整除 \(a\) 和 \(b\). 为证明 \(d\) 是最大公因子，设 \(c\) 是 \(a\) 与 \(b\) 的任意公因子，\(c = q_1^{g_1}q_2^{g_2}\cdots q_m^{g_m}\) 是 \(c\) 的素因子分解。由于每个 \(q_i\) 整除 \(c\)，因而整除 \(a\) 和 \(b\)，由前一命题我们看到 \(q_i\) 必整除某个 \(p_j\). 特别地（在相伴意义下，即乘以单位的意义下）\(c\) 中出现的素数必是 \(a\) 和 \(b\) 中出现的素数的子集：\(\{q_1, q_2, \ldots, q_m\} \subseteq \{p_1, p_2, \ldots, p_n\}\). 类似地，\(c\) 中出现的素数的指数必不大于 \(d\) 中出现的指数。这意味着 \(c\) 整除 \(d\)，完成证明。
\end{proof}

\begin{example}
在上面的例子中（\(a = 2210\)、\(b = 1131\)），我们由它们的素因子分解立即得到 \((a, b) = 13\). 注意若已知 \(a\) 和 \(b\) 的素因子分解，上面的命题立即给出它们的最大公因子，但求出这些素因子分解在计算上极其费时。欧几里得算法是确定两个整数最大公因子的最快方法，但遗憾的是它几乎不提供关于整数素因子分解的信息。
\end{example}

我们现在来到本章引入的一些环之间的一个主要结果。

\begin{theorem}\label{thm:8.14}
每个主理想整环都是唯一分解整环。特别地，每个欧几里得整环都是唯一分解整环。
\end{theorem}

\begin{proof}
注意第二个论断由第一个推出，因为欧几里得整环是主理想整环。为证明第一个论断，设 \(R\) 是主理想整环，\(r\) 是 \(R\) 的非零且不是单位的元素。我们必须首先证明 \(r\) 可以写成 \(R\) 的不可约元素的有限乘积，然后必须验证这个分解在单位的意义下唯一。

第一部分的证明方法与确定整数的素因子分解完全类似。假设 \(r\) 非零且不是单位。若 \(r\) 本身不可约，则完成。否则按定义 \(r\) 可以写成乘积 \(r = r_1r_2\)，其中 \(r_1, r_2\) 都不是单位。若这两个元素都不可约，则我们也完成，已经把 \(r\) 写成了不可约元素之积。否则这两个元素中至少有一个（设为 \(r_1\)）可约，因而可以写成两个非单位元素之积 \(r_1 = r_{11}r_{12}\)，如此继续。我们必须验证这个过程终止，即必在某个时刻达到使所得作为 \(r\) 的因子的所有元素都不可约。假设不然。由分解 \(r = r_1r_2\) 我们得到理想的真包含：\((r) \subsetneq (r_1) \subsetneq R\). 第一个包含是真的因为 \(r_2\) 不是单位，最后一个包含是真的因为 \(r_1\) 不是单位。由 \(r_1\) 的分解我们类似地得到 \((r) \subsetneq (r_1) \subsetneq (r_{11}) \subsetneq R\). 若这个分解过程不在有限步后终止，我们将得到理想的无限升链
\[
(r) \subsetneq (r_1) \subsetneq (r_{11}) \subsetneq \cdots \subsetneq R,
\]
其中所有包含都是真的，而选择公理保证无限链存在（参见附录 I）。

我们现在证明主理想整环中任何理想的升链 \(I_1 \subseteq I_2 \subseteq \cdots \subseteq R\) 最终稳定，即存在正整数 \(n\) 使对所有 \(k \geq n\) 有 \(I_k = I_n\). 特别地，不可能有所有包含都真的理想无限升链。设 \(I = \bigcup_{i=1}^{\infty} I_i\). 容易推出（如同 7.4 节命题 \ref{pro:7.11} 的证明）\(I\) 是理想。由于 \(R\) 是主理想整环，它是主的，设为 \(I = (a)\). 由于 \(I\) 是链中理想的并，\(a\) 必是链中某个理想的元素，设为 \(a \in I_n\). 但此时我们有 \(I_n \subseteq I = (a) \subseteq I_n\)，故 \(I = I_n\)，链在 \(I_n\) 处稳定。这证明 \(R\) 中每个非零且不是单位的元素都有某个分解为 \(R\) 中不可约元素的分解。

剩下要证明上面的分解本质唯一。我们对元素 \(r\) 的某个分解中不可约因子个数 \(n\) 作归纳。若 \(n = 0\)，则 \(r\) 是单位。如果我们有 \(r = qc\)（某个其他分解，\(q\) 不可约），则 \(q\) 将整除一个单位，因而它本身是单位，矛盾。现在假设 \(n \geq 1\) 且我们有两个乘积
\[
r = p_1p_2\cdots p_n = q_1q_2\cdots q_m
\]
是 \(r\) 的分解，其中 \(p_i, q_j\) 是（不必互异的）不可约元素。由于此时 \(p_1\) 整除右端的乘积，由命题 \ref{pro:8.11} 我们看到 \(p_1\) 必整除其中一个因子。必要时重新编号，可设 \(p_1\) 整除 \(q_1\). 但此时对 \(R\) 的某个元素 \(u\) 有 \(q_1 = p_1u\)，而由于 \(q_1\) 不可约，\(u\) 事实上必是单位。因此 \(p_1\) 和 \(q_1\) 相伴。消去 \(p_1\)（回忆我们在整环中，这是合法的），我们得到等式
\[
p_2\cdots p_n = q_2'\cdots q_m', \quad m \leq n,
\]
其中 \(q_2' = uq_2\) 仍不可约（与 \(q_2\) 相伴）。对 \(n\) 作归纳，我们得出结论：左端每个因子与右端（因而与中间，二者在相伴意义下相同）的因子一一对应（在相伴意义下）。由于 \(p_1\) 和 \(q_1\)（初始重新编号后）已经证明相伴，这完成了归纳步骤和定理的证明。
\end{proof}

\begin{corollary}[算术基本定理]\label{cor:8.15}
整数 \(\mathbb{Z}\) 是唯一分解整环。
\end{corollary}

\begin{proof}
整数 \(\mathbb{Z}\) 是欧几里得整环，故由定理是唯一分解整环。
\end{proof}

我们现在可以完成整环 \(R\) 上 Dedekind--Hasse 范数的存在性与 \(R\) 是否为 P.I.D. 之间的等价（命题 \ref{pro:8.9}）。

\begin{corollary}\label{cor:8.16}
设 \(R\) 是 P.I.D. 则 \(R\) 上存在乘性的 Dedekind--Hasse 范数。
\end{corollary}

\begin{proof}
若 \(R\) 是 P.I.D. 则 \(R\) 是 U.F.D. 定义范数 \(N\) 为 \(N(0) = 0\)、若 \(u\) 是单位则 \(N(u) = 1\)、若 \(a = p_1p_2\cdots p_n\)（\(p_i\) 是 \(R\) 中的不可约元素）则 \(N(a) = 2^n\)（良定义，因为 \(a\) 的不可约因子个数唯一）。显然 \(N(ab) = N(a)N(b)\)，故 \(N\) 是正的、乘性的。为证明 \(N\) 是 Dedekind--Hasse 范数，假设 \(a, b\) 是 \(R\) 的非零元素。则由假设由 \(a\) 和 \(b\) 生成的理想是主的，设为 \((a, b) = (r)\). 若 \(a\) 不含于理想 \((b)\) 中，则 \(r\) 也不含于 \((b)\) 中，即 \(r\) 不被 \(b\) 整除。由于对某个 \(x \in R\) 有 \(b = xr\)，可知 \(x\) 不是 \(R\) 中的单位，故 \(N(b) = N(x)N(r) > N(r)\). 因此 \((a, b)\) 含有范数严格小于 \(b\) 的范数的非零元素，完成证明。
\end{proof}

\subsection*{高斯整数中的分解}

我们以描述高斯整数 \(\mathbb{Z}[i]\) 中的不可约元素以及相应的对初等数论中 Fermat 的一个著名定理的应用来结束对唯一分解整环的讨论。这特别恰当，因为对 \(\mathbb{Z}[i]\) 的经典研究开创了环的代数研究。

一般地，设 \(\mathcal{O}\) 是二次整数环，\(N\) 是 7.1 节引入的相应域范数。假设 \(\alpha \in \mathcal{O}\) 的范数是 \(\mathbb{Z}\) 中的素数 \(p\). 若对某些 \(\beta, \gamma \in \mathcal{O}\) 有 \(\alpha = \beta\gamma\)，则 \(p = N(\alpha) = N(\beta)N(\gamma)\)，故 \(N(\beta)\) 或 \(N(\gamma)\) 之一是 \(\pm 1\)，另一个是 \(\pm p\). 由于我们已经看到 \(\mathcal{O}\) 的元素范数为 \(\pm 1\) 当且仅当它是 \(\mathcal{O}\) 中的单位，\(\alpha\) 的因子中有一个是单位。由此推出若 \(N(\alpha)\) 是 \(\pm\) 素数（在 \(\mathbb{Z}\) 中），则 \(\alpha\) 在 \(\mathcal{O}\) 中不可约。

假设 \(\pi\) 是 \(\mathcal{O}\) 中的素元，\((\pi)\) 是由 \(\pi\) 在 \(\mathcal{O}\) 中生成的理想。由于 \((\pi)\) 是 \(\mathcal{O}\) 中的素理想，容易验证 \((\pi) \cap \mathbb{Z}\) 是 \(\mathbb{Z}\) 中的素理想（若 \(a, b\) 是满足 \(ab \in (\pi)\) 的整数，则 \(a\) 或 \(b\) 是 \((\pi)\) 的元素，故 \(a\) 或 \(b\) 在 \((\pi) \cap \mathbb{Z}\) 中）。由于 \(N(\pi)\) 是 \((\pi)\) 中的非零整数，我们有 \((\pi) \cap \mathbb{Z} = p\mathbb{Z}\) 对某个整数素数 \(p\). 由此推出 \(p \in (\pi)\)，故 \(\pi\) 是整数素数 \(p\) 在 \(\mathcal{O}\) 中的因子，所以 \(\mathcal{O}\) 中的素元可以通过确定 \(\mathbb{Z}\) 中的素数如何在更大的环 \(\mathcal{O}\) 中分解来找到。假设 \(\pi\) 在 \(\mathcal{O}\) 中整除素数 \(p\)，设为 \(p = \pi\pi'\). 则 \(N(\pi)N(\pi') = N(p) = p^2\)，故由于 \(\pi\) 不是单位，只有两种可能：或者 \(N(\pi) = \pm p\)，或者 \(N(\pi) = \pm p^2\). 在前一情形 \(N(\pi') = \pm 1\)，故 \(\pi'\) 是单位，而 \(p = \pi\)（在相伴意义下）在 \(\mathbb{Z}[i]\) 中不可约。在后一情形 \(N(\pi) = N(\pi') = \pm p\)，故 \(\pi'\) 也不可约，而 \(p = \pi\pi'\) 恰好是两个不可约元素的乘积。

现在考虑高斯整数 \(\mathbb{Z}[i]\) 的特殊情形 \(D = -1\). 我们已经看到 \(\mathbb{Z}[i]\) 中的单位是 \(\pm 1\) 和 \(\pm i\). 我们在 8.1 节证明 \(\mathbb{Z}[i]\) 是欧几里得整环，因而是主理想整环和唯一分解整环，故不可约元素与素元相同，并可以通过看 \(\mathbb{Z}\) 中的素数如何在更大的环 \(\mathbb{Z}[i]\) 中分解来确定。

此时 \(\alpha = a+bi\) 有 \(N(\alpha) = \alpha\bar\alpha = a^2+b^2\)，其中 \(\bar\alpha = a-bi\) 是 \(\alpha\) 的复共轭。由刚才所看到的结果推出：\(p\) 在 \(\mathbb{Z}[i]\) 中恰好分解为两个不可约元素，当且仅当 \(p = a^2+b^2\) 是两个整数平方之和（否则 \(p\) 在 \(\mathbb{Z}[i]\) 中仍不可约）。若 \(p = a^2+b^2\)，则 \(\mathbb{Z}[i]\) 中相应的不可约元素是 \(a \pm bi\).

显然 \(2 = 1^2+1^2\) 是两个平方之和，给出分解 \(2 = (1+i)(1-i) = -i(1+i)^2\). 不可约元素 \(1+i\) 和 \(1-i = -i(1+i)\) 相伴，且容易验证这是共轭不可约元素 \(a+bi\) 与 \(a-bi\) 相伴的唯一情形。

由于任何整数的平方都同余于 0 或 1 模 4，\(\mathbb{Z}\) 中作为两个平方之和的奇素数必同余于 1 模 4. 因此若 \(p\) 是 \(\mathbb{Z}\) 中满足 \(p \equiv 3 \pmod 4\) 的素数，则 \(p\) 不是两个平方之和，而 \(p\) 在 \(\mathbb{Z}[i]\) 中仍不可约。

现在假设 \(p\) 是 \(\mathbb{Z}\) 中满足 \(p \equiv 1 \pmod 4\) 的素数。我们将证明 \(p\) 在 \(\mathbb{Z}[i]\) 中不可能不可约，这将表明 \(p = (a+bi)(a-bi)\) 分解为 \(\mathbb{Z}[i]\) 中两个不同不可约元素之积，等价地 \(p = a^2+b^2\) 是两个平方之和。我们首先证明初等数论中的如下结果：

\begin{lemma}\label{lem:8.17}
素数 \(p \in \mathbb{Z}\) 整除某个形如 \(n^2+1\) 的整数，当且仅当 \(p\) 是 2 或同余于 1 模 4 的奇素数。
\end{lemma}

\begin{proof}
\(p = 2\) 的陈述平凡，因为 \(2 \mid 1^2+1\). 若 \(p\) 是奇素数，注意 \(p \mid n^2+1\) 等价于 \(\mathbb{Z}/p\mathbb{Z}\) 中 \(n^2 = -1\). 这又等价于说 \(n\) 的剩余类在乘法群 \((\mathbb{Z}/p\mathbb{Z})^\times\) 中的阶为 4. 因此 \(p\) 整除某个形如 \(n^2+1\) 的整数当且仅当 \((\mathbb{Z}/p\mathbb{Z})^\times\) 含有 4 阶元素。由拉格朗日定理，若 \((\mathbb{Z}/p\mathbb{Z})^\times\) 含 4 阶元素，则 \(|(\mathbb{Z}/p\mathbb{Z})^\times| = p-1\) 被 4 整除，即 \(p\) 同余于 1 模 4.

反之，假设 \(p-1\) 被 4 整除。我们首先论证 \((\mathbb{Z}/p\mathbb{Z})^\times\) 含有唯一的 2 阶元素。若 \(m^2 \equiv 1 \pmod p\)，则 \(p\) 整除 \(m^2-1 = (m-1)(m+1)\). 因此 \(p\) 整除 \(m-1\)（即 \(m \equiv 1 \pmod p\)）或 \(m+1\)（即 \(m \equiv -1 \pmod p\)），故 \(-1\) 是 \((\mathbb{Z}/p\mathbb{Z})^\times\) 中唯一的 2 阶剩余类。现在阿贝尔群 \((\mathbb{Z}/p\mathbb{Z})^\times\) 含有 4 阶子群 \(H\)（例如，对子群 \(\{\pm 1\}\) 作商所得的群含有一个 2 阶子群，其原像是 \((\mathbb{Z}/p\mathbb{Z})^\times\) 中的 4 阶子群）。由于克莱因四元群有三个 2 阶元素，而 \((\mathbb{Z}/p\mathbb{Z})^\times\)（因而 \(H\)）有唯一的 2 阶元素，\(H\) 必是 4 阶循环群。因此 \((\mathbb{Z}/p\mathbb{Z})^\times\) 含有 4 阶元素，即 \(H\) 的一个生成元。
\end{proof}

\textbf{注：}我们将在后面证明 \((\mathbb{Z}/p\mathbb{Z})^\times\) 是循环群（9.5 节推论 19），由此立即可得存在 4 阶元素当且仅当 \(p-1\) 被 4 整除。

由引理 \ref{lem:8.17}，若 \(p \equiv 1 \pmod 4\) 是素数，则 \(p\) 在 \(\mathbb{Z}\) 中整除某个 \(n^2+1\)（\(n \in \mathbb{Z}\)），故当然在 \(\mathbb{Z}[i]\) 中 \(p\) 整除 \(n^2+1 = (n+i)(n-i)\). 若 \(p\) 在 \(\mathbb{Z}[i]\) 中不可约，则 \(p\) 将在 \(\mathbb{Z}[i]\) 中整除 \(n+i\) 或 \(n-i\). 在这种情形，由于 \(p\) 是实数，将推出 \(p\) 同时整除 \(n+i\) 及其复共轭 \(n-i\)；因而 \(p\) 将整除它们的差 \(2i\). 这显然不是事实。我们证明了如下结果：

\begin{proposition}\label{pro:8.18}
（1）（费马平方和定理）素数 \(p\) 是两个整数平方之和 \(p = a^2+b^2\)（\(a, b \in \mathbb{Z}\)）当且仅当 \(p = 2\) 或 \(p \equiv 1 \pmod 4\). 除交换 \(a, b\) 或改变 \(a, b\) 的符号外，\(p\) 表示为两个平方之和的方式唯一。

（2）高斯整数 \(\mathbb{Z}[i]\) 中的不可约元素如下：

（a）\(1+i\)（其范数为 2）；

（b）\(\mathbb{Z}\) 中满足 \(p \equiv 3 \pmod 4\) 的素数 \(p\)（其范数为 \(p^2\)）；

（c）\(a+bi\)、\(a-bi\)，即 \(\mathbb{Z}\) 中满足 \(p \equiv 1 \pmod 4\) 的素数 \(p = a^2+b^2 = (a+bi)(a-bi)\) 的互不相同的不可约因子（二者范数都是 \(p\)）。
\end{proposition}

命题 \ref{pro:8.18} 的第一部分是初等数论中 Fermat 的一个著名定理，对它可以给出若干不同的证明。

更一般地，整数 \(n \in \mathbb{Z}\) 能否写成两个整数平方之和 \(n = A^2+B^2\) 的问题等价于 \(n\) 是否为高斯整数中某个元素 \(A+Bi\) 的范数的问题，即 \(n = A^2+B^2 = N(A+Bi)\). 把 \(A+Bi = \pi_1\pi_2\cdots\pi_k\) 写成不可约元素之积（在单位意义下唯一），由命题 \ref{pro:8.18} 中 \(\mathbb{Z}[i]\) 中不可约元素的显式描述推出：\(n\) 是范数当且仅当 \(n\) 的素因子中同余于 3 模 4 的那些出现的指数为偶数。此外，若 \(n\) 的这个条件满足，则 \(A+Bi\) 在 \(\mathbb{Z}[i]\) 中分解的唯一性使我们能计数把 \(n\) 表示为两个平方之和的方式数，如下一推论所示。

\begin{corollary}\label{cor:8.19}
设 \(n\) 是正整数，写成
\[
n = 2^k p_1^{a_1}\cdots p_r^{a_r} q_1^{b_1}\cdots q_s^{b_s},
\]
其中 \(p_1, \ldots, p_r\) 是同余于 1 模 4 的互不相同素数，\(q_1, \ldots, q_s\) 是同余于 3 模 4 的互不相同素数。则 \(n\) 在 \(\mathbb{Z}\) 中可以写成两个平方之和，即 \(n = A^2+B^2\)（\(A, B \in \mathbb{Z}\)），当且仅当每个 \(b_i\) 都是偶数。此外，若 \(n\) 的这个条件满足，则把 \(n\) 表示为两个平方之和的方式数是
\[
4(a_1+1)(a_2+1)\cdots(a_r+1).
\]
\end{corollary}

\begin{proof}
推论的第一个陈述已在上面证明。现在假设 \(b_1, \ldots, b_s\) 都是偶数。对每个同余于 1 模 4 的素数 \(p_i\)，写 \(p_i = \pi_i\pi_i'\)（\(i = 1, 2, \ldots, r\)），其中 \(\pi_i\) 和 \(\pi_i'\) 是命题 \ref{pro:8.18}(2)(c) 中所述的不可约元素。若 \(N(A+Bi) = n\)，则检验范数我们看到（在单位意义下）\(A+Bi\) 在 \(\mathbb{Z}[i]\) 中的不可约分解为
\[
A+Bi = (1+i)^k(\pi_1^{a_{1,1}}\pi_1'^{a_{1,2}})\cdots(\pi_r^{a_{r,1}}\pi_r'^{a_{r,2}})q_1^{b_1/2}\cdots q_s^{b_s/2},
\]
其中非负整数 \(a_{i,1}, a_{i,2}\) 满足 \(a_{i,1}+a_{i,2} = a_i\)（\(i = 1, 2, \ldots, r\)）。因此 \(\mathbb{Z}[i]\) 中有 \((a_1+1)(a_2+1)\cdots(a_r+1)\) 个不同的范数为 \(n\) 的元素 \(A+Bi\)（在单位意义下）。最后，由于 \(\mathbb{Z}[i]\) 中有四个单位，推论的第二个陈述随之成立。
\end{proof}

\begin{example}
由于 \(493 = 17 \cdot 29\) 且两个素数都同余于 1 模 4，\(493 = A^2+B^2\) 是两个整数平方之和。由于 \(17 = (4+i)(4-i)\)、\(29 = (5+2i)(5-2i)\)，\(A+Bi\) 在 \(\mathbb{Z}[i]\) 中（在单位意义下）可能的分解是 \((4+i)(5+2i) = 18+13i\)、\((4+i)(5-2i) = 22-3i\)、\((4-i)(5+2i) = 22+3i\)、\((4-i)(5-2i) = 18-13i\). 乘以 \(-1\) 使两个符号都变号，乘以 \(i\) 交换 \(A\) 和 \(B\) 并引入一个符号变化。于是 \(493 = (\pm 18)^2+(\pm 13)^2 = (\pm 22)^2+(\pm 3)^2\)，所有可能的符号选择给出 493 表示为两个平方之和的 16 种可能方式中的 8 种；其余 8 种通过交换两个被加数得到。

类似地，整数 \(58000957 = 7^6 \cdot 17 \cdot 29\) 恰好可以以 16 种方式写成两个平方之和，这通过把上面 \(493 = A^2+B^2\) 中的整数 \(A, B\) 各乘以 \(343\) 得到。
\end{example}

\subsection*{总结}

总之，我们在带单位元的交换环的类之间有如下包含关系：
\[
\text{域} \subsetneq \text{欧几里得整环} \subsetneq \text{P.I.D.} \subsetneq \text{U.F.D.} \subsetneq \text{整环},
\]
且所有包含都是真包含。回忆 \(\mathbb{Z}\) 是欧几里得整环但不是域，二次整数环 \(\mathbb{Z}[(1+\sqrt{-19})/2]\) 是主理想整环但不是欧几里得整环，\(\mathbb{Z}[x]\) 是唯一分解整环（第 9 章定理 7）但不是主理想整环，而 \(\mathbb{Z}[\sqrt{-5}]\) 是整环但不是唯一分解整环。

\subsection*{习题}

\begin{enumerate}
\item 设 \(G = \mathbb{Q}^\times\) 是非零有理数的乘法群。若 \(a = \frac{p}{q} \in G\)，其中 \(p, q\) 是互素整数，设 \(\varphi : G \to G\) 是交换 \(p\) 和 \(q\) 的素幂分解中素数 2 和 3 的映射（例如 \(\varphi(2^4 3^1 5^1 13^2) = 3^4 2^1 5^1 13^2\)，\(\varphi(3^1 2^4) = 2^1 3^4\)，且 \(\varphi\) 在分子分母都与 2 和 3 互素的所有有理数上是恒等映射）。

（a）证明 \(\varphi\) 是群同构。

（b）证明群 \(G\) 到自身的同构有无穷多个。

（c）证明上面这些同构都不能延拓为环 \(\mathbb{Q}\) 到自身的同构。事实上证明恒等映射是 \(\mathbb{Q}\) 的唯一环同构。

\item 设 \(a\) 和 \(b\) 是唯一分解整环 \(R\) 的非零元素。证明 \(a\) 和 \(b\) 有最小公倍（参见第 1 节习题 11），并按命题 \ref{pro:8.13} 描述其最大公因子的同样方式，用 \(a\) 和 \(b\) 的素因子分解描述它。

\item 确定整数 \(2130797 = 17^2 \cdot 73 \cdot 101\) 表示为两个平方之和的所有方式。

\item 证明：若一个整数是两个有理数平方之和，则它是两个整数平方之和（例如 \(13 = (1/5)^2 + (18/5)^2 = 2^2+3^2\)）。

\item 设 \(R = \mathbb{Z}[\sqrt{-n}]\)，其中 \(n\) 是大于 3 的无平方因子整数。

（a）证明 \(2\)、\(\sqrt{-n}\) 和 \(1+\sqrt{-n}\) 在 \(R\) 中不可约。

（b）证明 \(R\) 不是 U.F.D. 由此推出二次整数环 \(\mathcal{O}\) 对 \(D \equiv 2, 3 \pmod 4\)、\(D < -3\) 不是 U.F.D.（因而也不欧几里得、不是 P.I.D.）。〔证明 \(\sqrt{-n}\) 或 \(1+\sqrt{-n}\) 之一不是素元。〕

（c）给出 \(R\) 中一个不是主理想的显式理想。〔用（b）考虑包含非素理想 \((\sqrt{-n})\) 或 \((1+\sqrt{-n})\) 的极大理想。〕

\item （a）证明商环 \(\mathbb{Z}[i]/(1+i)\) 是 2 阶域。

（b）设 \(q \in \mathbb{Z}\) 是满足 \(q \equiv 3 \pmod 4\) 的素数。证明商环 \(\mathbb{Z}[i]/(q)\) 是有 \(q^2\) 个元素的域。

（c）设 \(p \in \mathbb{Z}\) 是满足 \(p \equiv 1 \pmod 4\) 的素数，按命题 \ref{pro:8.18} 写 \(p = \pi\pi'\). 证明中国剩余定理（7.6 节定理 17）的假设满足，且作为环 \(\mathbb{Z}[i]/(p) \cong \mathbb{Z}[i]/(\pi) \times \mathbb{Z}[i]/(\pi')\). 证明商环 \(\mathbb{Z}[i]/(p)\) 的阶为 \(p^2\)，并由此推出 \(\mathbb{Z}[i]/(\pi)\) 和 \(\mathbb{Z}[i]/(\pi')\) 都是 \(p\) 阶域。

\item 设 \(\pi\) 是 \(\mathbb{Z}[i]\) 中的不可约元素。

（a）对任意整数 \(n \geq 0\)，证明 \((\pi^n)/(\pi^{n+1})\) 作为 \(\mathbb{Z}[i]\)-模是单模，且 \((\pi^n)/(\pi^{n+1}) \cong \mathbb{Z}[i]/(\pi)\) 作为 \(\mathbb{Z}[i]\)-模。

（b）证明 \(|\mathbb{Z}[i]/(\pi^n)| = |\mathbb{Z}[i]/(\pi)|^n\).

（c）对 \(\mathbb{Z}[i]\) 中任意非零 \(\alpha\)，证明商环 \(\mathbb{Z}[i]/(\alpha)\) 的阶等于 \(N(\alpha)\).〔用（b）连同中国剩余定理和前一习题的结果。〕

\item 设 \(R\) 是二次整数环 \(\mathbb{Z}[\sqrt{-5}]\)，定义理想 \(I_1 = (2, 1+\sqrt{-5})\)、\(I_2 = (3, 2+\sqrt{-5})\) 和 \(I_3 = (3, 2-\sqrt{-5})\).

（a）证明 \(2\)、\(3\)、\(1+\sqrt{-5}\) 和 \(1-\sqrt{-5}\) 在 \(R\) 中不可约、其中任意两个在 \(R\) 中不相伴，且 \(6 = 2 \cdot 3 = (1+\sqrt{-5})\cdot(1-\sqrt{-5})\) 是 6 在 \(R\) 中的两个不同不可约分解。

（b）证明 \(I_1\)、\(I_2\) 和 \(I_3\) 是 \(R\) 中的素理想。〔一种途径：对 \(I_3\) 注意由环的第三同构定理 \(R/I_3 \cong (R/(3))/(I_3/(3))\). 证明 \(R/(3)\) 有 9 个元素、\(I_3/(3)\) 有 3 个元素，且作为加法阿贝尔群 \(R/I_3 \cong \mathbb{Z}/3\mathbb{Z}\). 由此推出 \(I_3\) 是极大（因而素）理想，且作为环 \(R/I_3 \cong \mathbb{Z}/3\mathbb{Z}\).〕

（c）证明（a）中的分解蕴含理想等式 \((6) = (2)(3)\) 和 \((6) = (1+\sqrt{-5})(1-\sqrt{-5})\). 证明这两个理想分解给出理想 (6) 作为素理想之积的同一分解（参见第 2 节习题 5）。

\item 假设二次整数环 \(\mathcal{O}\) 是 P.I.D. 证明 \(\mathcal{O}\) 上域范数 \(N\) 的绝对值（参见 7.1 节）是 \(\mathcal{O}\) 上的 Dedekind--Hasse 范数。由此推出：若二次整数环 \(\mathcal{O}\) 拥有任何 Dedekind--Hasse 范数，则事实上 \(\mathcal{O}\) 上域范数的绝对值已经提供 \(\mathcal{O}\) 上的一个 Dedekind--Hasse 范数。〔若 \(\alpha, \beta \in \mathcal{O}\) 则对某个 \(\gamma \in \mathcal{O}\) 有 \((\alpha, \beta) = (\gamma)\). 证明若 \(\beta\) 不整除 \(\alpha\) 则 \(0 < |N(\gamma)| < |N(\beta)|\)——利用 \(\mathcal{O}\) 中的单位恰是范数为 \(\pm 1\) 的元素这一事实。〕

\noindent\textbf{注：}若 \(\mathcal{O}\) 关于某个范数是欧几里得整环，则它关于域范数的绝对值不必是欧几里得整环（尽管对 \(D < 0\) 这是对的，参见第 1 节习题 8）。一个例子是 \(D = 69\)（参见 D. Clark, A quadratic field which is Euclidean but not norm-Euclidean, Manuscripta Math., 83(1994), pp. 327--330）。

\item （\textbf{\(k\)-级欧几里得整环}）设 \(R\) 是整环，\(N : R \to \mathbb{Z}^+ \cup \{0\}\) 是 \(R\) 上的范数。环 \(R\) 关于 \(N\) 欧几里得，若对任意满足 \(b \neq 0\) 的 \(a, b \in R\)，存在 \(R\) 中的元素 \(q, r\) 使
\[
a = qb + r, \quad \text{其中 } r = 0 \text{ 或 } N(r) < N(b).
\]
现在假设这个条件被弱化，即对任意满足 \(b \neq 0\) 的 \(a, b \in R\)，存在 \(R\) 中的元素 \(q, q', r, r'\) 使
\[
a = qb + r, \quad b = q'r + r', \quad \text{其中 } r' = 0 \text{ 或 } N(r') < N(b),
\]
即两次除法之后的余数更小。称这样的整环为 \textbf{2-级欧几里得整环}。

（a）证明在 2-级欧几里得整环中迭代这些除法产生 \(a\) 与 \(b\) 的一个最大公因子，且它是 \(a\) 和 \(b\) 的线性组合。由此推出 2-级欧几里得整环的每个有限生成理想都是主理想。（然而存在不是 P.I.D. 的 2-级欧几里得整环。）〔模仿定理 \ref{thm:8.4} 的证明。〕

（b）证明每个非零非单位元素都可以分解为有限多个不可约元素的 2-级欧几里得整环是唯一分解整环。〔首先证明不可约元素是素元，如下。若 \(p\) 不可约且 \(p \mid ab\) 而 \(p\) 不整除 \(a\)，用（a）对某些 \(x, y\) 写 \(px + ay = 1\). 两边乘以 \(b\) 得出 \(p \mid b\)，故 \(p\) 是素元。然后沿用定理 \ref{thm:8.14} 中唯一性的证明。〕

（c）对任意整数 \(k \geq 1\) 作出明显的推广，定义 \textbf{\(k\)-级欧几里得整环}的概念。证明把“2-级欧几里得”换成“\(k\)-级欧几里得”后陈述（a）和（b）仍然成立。

\noindent\textbf{注：}存在是 2-级欧几里得但不是欧几里得的环的例子。也存在关于给定范数不欧几里得、但关于该范数是 \(k\)-级欧几里得的环的例子（例如，环 \(\mathbb{Z}[\sqrt{14}]\) 关于通常范数 \(N(a+b\sqrt{14}) = |a^2-14b^2|\) 不欧几里得，但关于这个范数是 2-级欧几里得的）。\(k\)-级欧几里得条件还与“元素取自 \(R\) 的可逆 \(n \times n\) 矩阵群 \(GL_n(R)\) 是否由初等矩阵（主对角线上全为 1、主对角线外恰有一个 1、其余全为 0 的矩阵）生成”这一问题有关。

\item （\textbf{P.I.D. 的刻画}）证明 \(R\) 是 P.I.D. 当且仅当 \(R\) 是同时也是 Bézout 整环的 U.F.D.（参见第 2 节习题 7）。〔一个方向由定理 \ref{thm:8.14} 给出。反之，设 \(a\) 是理想 \(I\) 中具有最少不可约因子个数的非零元素。通过证明若存在 \(b \in I\) 而 \(b \notin (a)\) 则 \((a, b) = (d)\) 导致矛盾，来证明 \(I = (a)\).〕
\end{enumerate}
